The following parts of the survey are structured as such: the survey starts with a comprehensive review of the essential aspects of MLLMs, including 
(1) Mainstream architectures (\S \ref{sec:mllm_arch});
(2) A full recipe of training strategy and data (\S \ref{sec:mllm_train}); (3) Common practices of performance evaluation (\S \ref{sec:mllm_eval}). 
Then, we delve into a deeper discussion on some important topics about MLLMs, each focusing on a main problem: (1) What aspects can be further improved or extended (\S \ref{sec:extensions})? (2) How to relieve the multimodal hallucination issue (\S \ref{sec:hallu})?
The survey continues with the introduction of three key techniques (\S \ref{sec:tech}), each specialized in a specific scenario: M-ICL (\S \ref{sec:micl}) is an effective technique commonly used at the inference stage to boost few-shot performance. Another important technique is M-CoT (\S \ref{sec:mcot}), which is typically used in complex reasoning tasks. Afterward, we delineate a general idea to develop LLM-based systems to solve composite reasoning tasks or to address common user queries (\S \ref{sec:vr}).
Finally, we finish our survey with a summary and potential research directions.







